\clearpage
\begin{trapbox}
\textbf{误区一：只要两个边切圆，就认为第三边也切圆}

\textbf{错法：}圆心在 $x$ 轴上就够了，随便选 $h$。

\textbf{为什么错：}见例3，$C:y^2=x$、$r=1$，当 $h=3$ 时，两条过 $O$ 的切线再次交抛物线于 $(8,\pm2\sqrt2)$，第三条直线 $x=8$ 距圆心为 $5$，不相切。

\textbf{怎样改：}先核对特定位置 $h=r+r^2/k$。另一个分支 $h=r^2/k-r$ 也能使三\emph{直线}相切，但不能因此认为圆就是三角形内切圆。

\medskip
\textbf{误区二：用斜率写弦，悄悄丢掉竖直线}

\textbf{错法：}把 $kx-(u+v)y+uv=0$ 改为 $y=\frac{k}{u+v}x+\frac{uv}{u+v}$，不检查分母。

\textbf{怎样改：}保留 $kx-(u+v)y+uv=0$。当 $u+v=0$ 时，它变为 $x=u^2/k$，恰好是竖直弦，并非没有直线。
\end{trapbox}

\clearpage
\begin{trapbox}
\textbf{误区三：把“直线相切”说成“每段边上都有切点”}

\textbf{反例（精确坐标）：}取 $C:y^2=x$、$\Gamma:x^2+y^2=1$，对应另一分支 $h=r^2/k-r=0$。固定
$A_1=(9/4,3/2)$，由两条切圆的弦得到另两点的纵坐标
\[
b=\frac{6+\sqrt{101}}5,\qquad c=\frac{6-\sqrt{101}}5.
\]
它们的和为 $12/5$、积为 $-13/5$，因此第三条弦
\[
A_2A_3:\ 5x-12y-13=0
\]
到圆心的距离是 $13/13=1$，但切点（圆心到直线的垂足）是
\[
T=(5/13,-12/13)\approx(0.385,-0.923).
\]
而 $A_3=(c^2,c)\approx(0.656,-0.810)$。因此 $T$ 位于 $A_3$ 朝着\textbf{远离 $A_2$}的延长线上，不属于线段 $A_2A_3$。
\begin{center}
% 第二分支：固定 a=3/2, b,c=(6±sqrt101)/5；第三弦5x-12y-13=0。
% 这里只画靠近A3的局部，突出A3外侧的切点，未画到远处的A2。
\begin{tikzpicture}[x=1.65cm,y=1.65cm,>=Stealth,font=\small]
\pgfmathsetmacro{\ccc}{(6-sqrt(101))/5}
\draw[gray!50,->] (-1.15,0)--(3.05,0) node[right] {$x$};
\draw[gray!50,->] (0,-1.45)--(0,1.50) node[above] {$y$};
\draw[blue!70!black,thick,domain=-1.22:1.22,samples=110,variable=\t] plot ({\t*\t},{\t});
\draw[teal!80!black,thick,fill=teal!8] (0,0) circle(1);
\coordinate (C) at ({\ccc*\ccc},\ccc);
\coordinate (T) at ({5/13},{-12/13});
\coordinate (E) at (2.8, { (5*2.8-13)/12 });
\draw[orange!85!black,very thick,->] (C)--(E) node[above right,black,font=\scriptsize] {沿线段趋向 $A_2$};
\draw[red!75!black,very thick,densely dashed] (T)--(C);
\draw[teal!80!black,densely dashed] (0,0)--(T);
\fill[red!75!black] (T) circle (1.8pt);
\fill[orange!85!black] (C) circle (1.8pt);
\node[below left,red!75!black] at (T) {$T$};
\node[right] at (1.15,-.94) {$A_3$};
\node[above left] at (0,0) {$M=O$};
\node[font=\small,red!75!black] at (1.2,-1.72) {红虚线：切点在边的延长线};
\end{tikzpicture}

\end{center}
\textbf{怎样改：}问直线与圆相切，只检查圆心到直线距离等于半径；问三角形内切圆，必须继续检查切点在三条线段内部。
\end{trapbox}

\clearpage
\begin{trapbox}
\textbf{误区四：相减求出了和，却仍然把两个未知数逐个求出来}

\textbf{错法：}先求 $b+c$，然后去解二次方程分别得到根式 $b,c$，再计算 $bc$。

\textbf{为什么绕远：}目标直线 $kx-(b+c)y+bc=0$ 只要和与积。两式相减得到 $b+c$，相加后利用
\[
b^2+c^2=(b+c)^2-2bc
\]
立刻能求 $bc$。不需要认出同构，也不需要把两个未知数分别解出。

\medskip
\textbf{误区五：随便约去 $b-c$ 或 $a^2-r^2$}

\textbf{怎样改：}三点不同保证 $b-c\ne0$。若 $a^2=r^2$，则相减式要求 $0=-2kra\ne0$，与题设矛盾；\textbf{说明合法情形中 $a^2-r^2\ne0$，不是能不加解释地除。}

\medskip
\textbf{误区六：把“同一条件”误当成“同一根”}

\textbf{错法：}两条切线条件相同就写 $b=c$。

\textbf{怎样改：}它们是\textbf{同一个二次方程的两个不同解}；正是 $b\ne c$ 才能从差式中约去 $b-c$。第三条弦仍由 $b+c$ 和 $bc$ 决定。
\end{trapbox}
